\section{Comparison with the actual flow}
\label{sec:actual-comparison}

Assume that \(W\) is reduced and \(L\not\simeq M\).
Let \(H(T)\) solve \eqref{eq:metric-flow} from any smooth
initial metric with determinant \(h_{\det}\), and let
\(\overline H(T)=\widehat H_{a(T)}\) be the curve in
Proposition~\ref{prop:comparison-trajectory}. No smallness
condition is imposed on the initial distance between them.

\begin{theorem}\label{thm:actual-comparison}
There are \(C>0\) and \(0<\gamma<1\), depending on the
fixed data and the initial metric, such that
\begin{equation}\label{eq:distance-decay}
 \sup_X\left|\log\bigl(\overline H(T)^{-1}H(T)\bigr)\right|_{\overline H(T)}
 \le C(1+T)^{-\gamma}.
\end{equation}
Consequently
\begin{equation}\label{eq:actual-mass}
 Z(T)\sim\frac{v_0^2}{4\pi^2\sigma T}.
\end{equation}
With the intrinsic amplitude \(\alpha(T)=Z(T)/v_0\),
\begin{equation}\label{eq:global-inclusion-profile}
 \sup_X\left|\frac{u(T)}{Z(T)}
       -\frac1{v_0}\frac{u_0}{u_0+\alpha(T)}\right|
 \le C(1+T)^{-\gamma}.
\end{equation}
The same estimate holds with \(a(T)\) in place of \(\alpha(T)\).
There are fixed \(r_0,T_0>0\) such that for \(0<r\le r_0\)
and \(T\ge T_0\),
\begin{equation}\label{eq:shrinking-exterior}
 \sup_{X\setminus B_r(\Sigma)}
 \left|\frac{u(T)}{Z(T)}-\frac1{v_0}\right|
 \le C\left(T^{-\gamma}+\frac1{Tr^2}\right).
\end{equation}
\end{theorem}

\begin{proof}
At each point the metrics with fixed determinant form the
positive-matrix symmetric space with tangent norm
\(\tr((H^{-1}\delta H)^2)\). Put
\[
 s=\log(\overline H^{-1}H),\qquad
 H_\theta=\overline He^{\theta s},\qquad
 d=|s|_{\overline H},\qquad e=\tfrac12d^2.
\]
The endomorphism \(s\) is trace-free and Hermitian for every
\(H_\theta\). The metric-distance calculation for two
solutions of HYM flow, in the convention
\eqref{eq:metric-flow}, is
\[
 (\partial_T+\Delta_0)e
 =-2\int_0^1|D'_{H_\theta}s|_{H_\theta}^2\,\dd\theta.
\]
This is \cite[Lemma~6.4.5]{Kobayashi87}, with book time
\(t=2T\) and scalar operator \(\Delta_0/2\).
For our second endpoint the additional velocity is
\(\mathcal R=\overline H^{-1}\overline H_T+2\Phi_{\overline H}\).
The first variation of squared symmetric-space distance
has endpoint gradient of norm \(d\). Thus, with
\(\varepsilon(T)=\|\mathcal R(T)\|_\infty\),
\begin{equation}\label{eq:forced-square-distance}
 (\partial_T+\Delta_0)e
 \le-2\int_0^1|D'_{H_\theta}s|_{H_\theta}^2\,\dd\theta
       +\varepsilon(T)d.
\end{equation}
This uses only the endpoint velocity, with no derivative
of \(\mathcal R\).

For every \(\theta\), metric compatibility and Kato's
inequality give
\[
 |\nabla d|^2\le|D_{H_\theta}s|^2
              =2|D'_{H_\theta}s|^2.
\]
Average over \(\theta\) and apply the chain rule to
\eqref{eq:forced-square-distance}. The derivative term
has the sign needed to give
\begin{equation}\label{eq:forced-distance}
 (\partial_T+\Delta_0)d\le\varepsilon(T).
\end{equation}
At \(d=0\) use \((d^2+\eta^2)^{1/2}\) and then
\(\eta\downarrow0\). The maximum principle and
\eqref{eq:integrable-flow-error} imply
\begin{equation}\label{eq:bounded-distance}
 \sup_Xd(T)\le
 B:=\sup_Xd(0)+\int_0^\infty\varepsilon(t)\,\dd t<\infty.
\end{equation}
Every smooth positive initial metric on the compact
bundle gives a finite \(B\).

We use the first eigenvalue in
\eqref{eq:sharp-small-eigenvalue} to improve boundedness
to decay. Choose \(a_{\mathrm{in}}\) small enough that
\(\lambda_1(\widehat{\mathcal L}_a)\ge c_1a\) throughout
the comparison trajectory. The tensor \(D''s\) is the
same for all \(H_\theta\), because the holomorphic
structure is fixed. Its squared fibre norms differ
by factors between \(e^{-2B}\) and \(e^{2B}\).
For example, after diagonalizing \(s\) at a point,
a matrix entry has factor \(e^{\theta(s_i-s_j)}\).
As the \(D'\) and \(D''\) norms of a Hermitian
section agree, it follows that
\[
 \int_X|D'_{H_\theta}s|^2\,\dd V
 \ge e^{-2B}\widehat{\mathfrak q}_{a(T)}(s,s)
 \ge e^{-2B}c_1a(T)\int_Xd^2\,\dd V.
\]
Put \(D(T)=\|d(T)\|_2\). Integration of
\eqref{eq:forced-square-distance} and Cauchy--Schwarz
give, with the usual regularization at \(D=0\),
\begin{equation}\label{eq:distance-ode}
 D'(T)\le-2e^{-2B}c_1a(T)D(T)+\sqrt{v_0}\,\varepsilon(T).
\end{equation}
Here \(a(T)\asymp(1+T)^{-1}\), while
\[
 \varepsilon(T)\le
 C(1+T)^{-2}(1+\log(2+T))^3.
\]
Choose \(0<\gamma<1\) smaller than an eventual
coefficient of \((1+T)^{-1}\) in the damping term.
Multiplying \eqref{eq:distance-ode} by \((1+T)^\gamma\)
and integrating gives \(D(T)\le C(1+T)^{-\gamma}\),
since
\(\int_1^\infty t^{\gamma-2}(1+\log t)^3\,\dd t<\infty\).
The fixed scalar heat comparison for
\eqref{eq:forced-distance} on \([T-1,T]\) now gives
\[
 \|d(T)\|_\infty
 \le C D(T-1)+\int_{T-1}^T\varepsilon(t)\,\dd t
 \le C(1+T)^{-\gamma}.
\]
This proves \eqref{eq:distance-decay}.

The correction from \(H_{a(T)}\) to \(\overline H(T)\)
has supremum logarithm \(O(a\ell(a))\). Decreasing
\(\gamma\) if needed, the preceding estimate yields
the relative quadratic-form bounds
\begin{equation}\label{eq:actual-static-comparison}
 e^{-CT^{-\gamma}}H_{a(T)}
 \le H(T)\le e^{CT^{-\gamma}}H_{a(T)}.
\end{equation}
Evaluation on the fixed inclusion and integration give
\[
 u(T,x)=u_{a(T)}(x)(1+O(T^{-\gamma})),\qquad
 Z(T)=Z_{a(T)}(1+O(T^{-\gamma})),
\]
with uniform errors off \(\Sigma\); both inclusion
norms vanish on \(\Sigma\).
For reduced \(W\), \eqref{eq:balanced-static} and
\eqref{eq:compact-identities} say
\[
 Z_a=av_0+O(a^2),\qquad u_a=\frac{au_0}{u_0+a}.
\]
Together with \eqref{eq:comparison-asymptotic}, these
prove \eqref{eq:actual-mass} and the profile with
parameter \(a(T)\). They also give
\(\alpha(T)/a(T)=1+O(T^{-\gamma})\).
For \(u_0\ge0\), changing \(a\) to \(\alpha\) in
\(u_0/(u_0+a)\) changes the result by at most
\(C|\alpha/a-1|\), when this ratio is bounded above
and below. This proves \eqref{eq:global-inclusion-profile}.

Near the reduced zeros \(u_0\asymp r^2\), and \(u_0\)
is bounded below on the fixed exterior. Hence
\(u_0\ge cr^2\) on \(X\setminus B_r(\Sigma)\).
The error in replacing \(u_0/(u_0+a)\) by one is
at most \(a/u_0\le C/(Tr^2)\), proving
\eqref{eq:shrinking-exterior}.
\end{proof}

In particular, a sufficient time for an exterior error at most
\(\epsilon\), with constants depending only on the specified
data and initial metric, is
\[
 T\ge C\max\{1,\epsilon^{-1/\gamma},(\epsilon r^2)^{-1}\}.
\]
This conclusion has joint time and radius quantifiers. It
follows from the global comparison, rather than from qualitative
convergence on a fixed compact set.

\subsection{The intrinsic scale of the inclusion}

For each reduced point put
\begin{equation}\label{eq:intrinsic-radius}
 r_p(T)=\sqrt{\frac{\alpha(T)}{\lambda_p}}.
\end{equation}
In the fixed normalized holomorphic coordinates and frames,
define
\begin{equation}\label{eq:inner-metric-limit}
 H_{\mathrm{loc}}(y)
 =(1+|y|^2)I-\frac{|y|^2+2}{1+|y|^2}
      (y_1,-y_2)^{\mathsf t}(\bar y_1,-\bar y_2).
\end{equation}
This is the metric of Lemma~\ref{lem:local-metric} at
parameter one, with the fixed unitary sign change.

\begin{corollary}\label{cor:c0-inner}
As \(T\to\infty\), uniformly on each compact \(y\)-set,
\begin{equation}\label{eq:c0-inner}
 H(T,r_p(T)y)\longrightarrow H_{\mathrm{loc}}(y),\qquad
 \frac{u(T,r_p(T)y)}{\alpha(T)}
       \longrightarrow\frac{|y|^2}{1+|y|^2}.
\end{equation}
Along each radial ray at \(p\), the first radius where
\(u(T,\cdot)=\alpha(T)/2\) is
\(r_p(T)(1+o(1))\), uniformly in direction.
\end{corollary}

\begin{proof}
The normalized coordinates give \(u_0(z)=c(z)|z|^2\)
with \(c(z)=\lambda_p+O(|z|^2)\). Thus
\(u_0(r_py)/\alpha\to|y|^2\) smoothly on compact
sets. Since \(\alpha/a\to1\), the exact compact
metric formula converges there to
\eqref{eq:inner-metric-limit}; its determinant weight
tends to one. The limit and its inverse are bounded
on each compact set. The relative estimate
\eqref{eq:actual-static-comparison} therefore proves
the first assertion in \eqref{eq:c0-inner}.
Evaluation on
\(\iota(r_py)/\sqrt\alpha=(y_1,-y_2)^{\mathsf t}\)
proves the second.
For \(0<\delta<1\), the scalar limit is strictly
below \(1/2\) on \(|y|\le1-\delta\), and strictly
above \(1/2\) on \(|y|=1+\delta\).
For all sufficiently large \(T\) the same statements
hold uniformly for the actual scalar norm. Every ray
has a crossing between these two radii, and none
before the smaller one. Letting \(\delta\downarrow0\)
proves the claimed first-crossing equivalent.
\end{proof}

The metric convergence just proved is in \(C^0\). It does not
yet imply convergence of curvature, and the radius defined by
the inclusion need not agree by definition with any curvature
radius. We establish the higher regularity and compare the two
radii in the next section.
